%%% bubblequiz — static half of the bubble sheet.
%%%
%%% R/make_bubblesheet.R writes output/bubblesheet_v<N>.tex, which \def's the
%%% macros used below, \input's this file, then emits the section bars and
%%% answer-row tabulars derived from exam.yml.
%%%
%%% Macros expected from the generated file:
%%%   \examversion \examtitle \examdate \examsubtitle \examdurationmarks
%%%   \examinstructions \idlabel \idprefix \idexample \idlastcol \idblanks
%%%
%%% DO NOT change the bubble geometry (8pt circles, 1.2pt lines, 1.8cm column
%%% gap, 4pt row skip) or the corner registration squares without re-running
%%% `make calibrate` — the marker locates bubbles from these positions.

\documentclass[11pt, a4paper]{article}
\usepackage[a4paper, left=1.4cm, right=1.4cm, top=0.9cm, bottom=0.9cm]{geometry}
\usepackage{fontspec}
\setmainfont{Helvetica Neue}
\usepackage{tikz}
\usepackage{xcolor}
\definecolor{shadedbg}{gray}{0.88}
\definecolor{versionbg}{gray}{0.15}

\providecommand{\examversion}{?}

\pagestyle{empty}
\setlength{\parskip}{0pt}
\setlength{\parindent}{0pt}

%── Answer bubble: 7 mm circle with letter centred ───────────────────────────
\newcommand{\bub}[1]{%
  \begin{tikzpicture}[baseline=-0.6ex]
    \draw[line width=1.2pt](0,0) circle (8pt);
    \node[font=\fontsize{8.5}{8.5}\selectfont\bfseries] at (0,0){#1};
  \end{tikzpicture}%
  \hspace{4pt}%
}

%── One question row: Q<n>  ○A ○B ○C ○D ○E ──────────────────────────────────
%   #1 = question number, #2 = option letters (emitted by make_bubblesheet.R)
\newcommand{\qrow}[2]{%
  \makebox[2.4em][l]{\textbf{\small Q#1}}%
  #2%
}

%── An essay placeholder row: Q<n>  Essay <marks> marks ─────────────────────
\newcommand{\erow}[2]{%
  \makebox[2.4em][l]{\textbf{\small Q#1}} Essay #2 marks%
}

%── Shaded section heading bar ───────────────────────────────────────────────
\newcommand{\sectionbar}[2]{%
  \vspace{4pt}%
  \noindent\colorbox{shadedbg}{\parbox{\dimexpr\linewidth-2\fboxsep}{%
    \vspace{1pt}\textbf{#1}\hfill{\small\textit{#2}}\vspace{1pt}%
  }}\par\vspace{3pt}%
}

\begin{document}

%── Registration marks: solid 8 mm squares in each corner (OMR anchors) ──────
%   Squares preferred over circles: sharp corners give precise skew correction.
\begin{tikzpicture}[remember picture, overlay]
  \fill[black] ([xshift= 3mm, yshift= -3mm]current page.north west) rectangle ++( 5mm, -5mm);
  \fill[black] ([xshift=-8mm, yshift= -3mm]current page.north east) rectangle ++( 5mm, -5mm);
  \fill[black] ([xshift= 3mm, yshift=  3mm]current page.south west) rectangle ++( 5mm,  5mm);
  \fill[black] ([xshift=-8mm, yshift=  3mm]current page.south east) rectangle ++( 5mm,  5mm);
\end{tikzpicture}

%── Title block ───────────────────────────────────────────────────────────────
\noindent{\LARGE\bfseries \examtitle}%
\hfill\colorbox{versionbg}{\textcolor{white}{\Large\bfseries\quad Version \examversion\quad}}
\hspace{4pt}{\large\bfseries \examdate}\\[2pt]
{\large \examsubtitle}%
\hfill{\normalsize \examdurationmarks}

\vspace{3pt}
\noindent\rule{\linewidth}{1.4pt}
\vspace{2pt}

%── Instructions ──────────────────────────────────────────────────────────────
{\small \examinstructions}

\vspace{5pt}

%── Student details (left) + ID digit grid (right), side by side ─────────────
\noindent
\begin{minipage}[t]{0.40\linewidth}
  \textbf{Name}\\[4pt]
  \underline{\hspace{\linewidth}}\\[12pt]
  \textbf{\idlabel} \quad {\small\textit{e.g.\ \idexample}}\\[5pt]
  \idprefix\,\idblanks\\[8pt]
  {\small Also {\bf fill your \idlabel\ in digit bubbles}, this enables automatic read  $\rightarrow$} \\[8pt]
  {\small {\bf Tick to confirm} the version number of the test (top right corner) is the same as your answer booklet}  \bub{ }\\[8pt]
\end{minipage}%
\hfill
\begin{minipage}[t]{0.56\linewidth}
  \textbf{\small \idlabel\ bubbles} {\small — one bubble per column, left to right:}\\[3pt]
  \begin{tikzpicture}[x=1.08cm, y=-0.70cm]
    \foreach \col in {0,...,\idlastcol} {
      \foreach \d in {0,...,9} {
        \draw[line width=0.9pt](\col, \d) circle (0.24cm);
        \node[font=\fontsize{7.5}{7.5}\selectfont] at (\col, \d) {\d};
      }
    }
  \end{tikzpicture}
\end{minipage}

\vspace{4pt}
\noindent\rule{\linewidth}{1.4pt}%
